% 张三慧第4版V1第7章《波动》独立计算候选解
% 原教材PDF物理页248--253；配套辅导书对照/原图重绘尚未验收。
\section*{第7章\quad 波动：独立解析}

\paragraph{习题 7.1}
频率等于行波速度除以波长。
题给
\(u=740\,\mathrm{km/h}\)，
\(\lambda=300\,\mathrm{km}\)，所以
\[
 \boxed{f=\frac u\lambda
 =\frac{740}{300\times3600}
 \approx6.85\times10^{-4}\,\mathrm{Hz}}.
\]
对应周期约 \(24.3\,\mathrm{min}\)。
横跨 \(L=8000\,\mathrm{km}\)
所需时间
\[
 \boxed{t=\frac Lu=\frac{8000}{740}
 \approx10.81\,\mathrm h}.
\]
按题给的恒速波传播模型，不计不同水深的色散。

\paragraph{习题 7.2}
在 \(x_0=0.10\,\mathrm m\) 处
\(y(x_0,t)=0.05\sin(1.0-4.0t)\)。
角频率 \(\omega=4\,\mathrm{rad/s}\)，
沿 \(+x\) 传播时
\(k=\omega/u=4/0.80=5.0\,\mathrm{rad/m}\)，
行波相位随 \((x-x_0)\) 增加 \(k(x-x_0)\)，
因此
\[
 \boxed{y(x,t)=0.050
 \sin[1.0-4.0t+5.0(x-0.10)]\ \mathrm m}
 =0.050\sin(0.50+5x-4t)\,\mathrm m.
\]
\textbf{方向校核：}注意
\(\sin(kx-\omega t)\) 确实向 \(+x\)
传播，而 \(\sin(kx+\omega t)\) 反向传播。

\paragraph{习题 7.3}
原波函数
\(y=0.020\sin[2\pi(200t-2x)]\)，
故
\[
 \boxed{f=200\,\mathrm{Hz},\quad
 \lambda=1/2=0.50\,\mathrm m,\quad
 u=f\lambda=100\,\mathrm{m/s}}.
\]
相位固定时 \(200t-2x=C\)，
因此波沿 \(+x\) 传播。
绳上质元最大振动速率
\(v_{\max}=A\omega
 =0.020(2\pi 200)\approx25.13\,\mathrm{m/s}\)，
与波速之比
\[
 \boxed{v_{\max}/u\approx0.2513}.
\]
质点振动速度不能和波形传播速度混同。

\paragraph{习题 7.4}
简谐地震波的
\(f=1\,\mathrm{Hz}\)、
波速 \(u=3000\,\mathrm{m/s}\)，
波长
\[
 \boxed{\lambda=\frac uf=3.00\times10^3\,\mathrm m}.
\]
地面垂直振幅 \(A=2.0\,\mathrm m\)，
最大速度
\[
 \boxed{|v|_{\max}=2\pi fA=4\pi
 \approx12.57\,\mathrm{m/s}}.
\]
这是题设高度理想化的单频地震波，
不是对实际地震时完整运动过程的预测。

\paragraph{习题 7.5}
原图7.44可读振幅
\(A=0.040\,\mathrm m\)、波长
\(\lambda=0.40\,\mathrm m\)，
沿 \(+x\) 传播速度
\(u=0.080\,\mathrm{m/s}\)，
故
\[
 k=\frac{2\pi}{\lambda}=5\pi\,\mathrm{rad/m},\quad
 \omega=uk=0.40\pi\,\mathrm{rad/s},\quad
 T=\frac{\lambda}{u}=5.00\,\mathrm s.
\]
由 \(t=0\) 原点向上过零的\emph{空间曲线}，
取
\[
 \boxed{y(x,t)=0.040\sin(5\pi x-0.40\pi t)\ \mathrm m}.
\]
注意波形沿 \(+x\) 运动，
而 \(x=0,t=0\) 处的介质质点速度
\(\partial_ty=-A\omega<0\)，
不能把波形空间上升斜率误当成质点向上速度。
当 \(t=T/8=0.625\,\mathrm s\)，
波形相当于初图右移
\(uT/8=\lambda/8=0.050\,\mathrm m\)：
\[
 \boxed{y(x,T/8)=0.040
 \sin(5\pi x-\pi/4)\ \mathrm m}.
\]
原图矢量重绘仍待复核。

\paragraph{习题 7.6}
给定
\(y=A\cos[\pi(4t+2x)]\)，
有 \(\omega=4\pi\)，
\(k=2\pi\)；
因为相位 \(4t+2x\) 恒定需
\(dx/dt=-2\,\mathrm{m/s}\)，
波向 \(-x\) 传播。
在 \(t=4.2\,\mathrm s\) 时
\[
 \boxed{y(x,4.2)=A\cos(0.8\pi+2\pi x)},
\]
其中相位已约去整周。
波峰满足 \(4t+2x=2n\)，
所以
\(x=n-8.4\,\mathrm m\)。
离原点最近的一峰在
\[
 \boxed{x_{\rm near}=-0.40\,\mathrm m}
\]
（右侧相邻峰位于 \(+0.60\,\mathrm m\)）。
波向左行进，所以右侧波峰从
\(x=0.60\) 到原点还需
\(0.60/2=0.30\,\mathrm s\)，
即经过原点的
\[
 \boxed{t=4.50\,\mathrm s}.
\]
此刻图形按上述余弦式直接绘出。

\paragraph{习题 7.7}
\(u=350\,\mathrm{m/s}\)，
\(f=500\,\mathrm{Hz}\)，
波长 \(\lambda=u/f=0.70\,\mathrm m\)。
若两点相位差为
\(60^\circ=\pi/3\)，
取最近的正距离：
\[
 \boxed{\Delta x=\frac{\Delta\varphi}{2\pi}\lambda
 =\frac{0.70}{6}\approx0.117\,\mathrm m}.
\]
同一质点相隔 \(10^{-3}\,\mathrm s\)：
\[
 \boxed{\Delta\varphi
 =2\pi f\Delta t
 =2\pi(500)(10^{-3})=\pi}.
\]
在相位差中可任选加整周 \(2n\pi\)，
但物理状态由模 \(2\pi\) 决定。

\paragraph{习题 7.8}
均匀细杆的一维纵波速率为
\(u=\sqrt{Y/\rho}\)，
其中 \(Y\) 是杨氏模量。
由题给参数
\[
 \boxed{Y=\rho u^2
 =(7.8\times10^3)(5100)^2
 \approx2.03\times10^{11}\,\mathrm{Pa}}.
\]
这是理想无限细均匀弹性杆的
无色散纵波近似。

\paragraph{习题 7.9}
体积弹性模量定义
\(K=-V\,dp/dV\)。
若 \(|\Delta V|\ll V\)，
则小形变的恢复压强增量大小
\(\Delta p\simeq K|\Delta V|/V\)。
从零应变到最终压缩量所作弹性储能功
\[
 W=\int_0^{|\Delta V|}
 \frac{K\,\delta V}{V}\,d(\delta V)
 =\frac{K(\Delta V)^2}{2V}.
\]
除以原体积 \(V\)，得
\[
 \boxed{u_{\rm elastic}
 =W/V=\frac12K(\Delta V/V)^2}.
\]
压缩与膨胀的能量均为正，原因是
它依赖形变的平方。

\paragraph{习题 7.10}
均匀一维简谐纵波平均强度
\(I_{\rm wave}=\tfrac12\rho u\omega^2A^2\)，
其中 \(\omega=2\pi f\)。
代入
\(\rho=7.9\times10^3\,\mathrm{kg/m^3}\)，
\(u=5.1\times10^3\,\mathrm{m/s}\)，
\(f=10^3\,\mathrm{Hz}\)，
\(A=10^{-9}\,\mathrm m\)，有
\[
 \boxed{I_{\rm wave}
 =\tfrac12(7900)(5100)(2\pi1000)^2(10^{-9})^2
 \approx7.95\times10^{-4}\,\mathrm{W/m^2}}.
\]
钢轨截面积
\(S=15\,\mathrm{cm^2}
 =1.5\times10^{-3}\,\mathrm{m^2}\)，
平均传输功率
\[
 \boxed{P=I_{\rm wave}S
 \approx1.19\times10^{-6}\,\mathrm W}.
\]
注意声强单位为 \(\mathrm{W/m^2}\)，
与声级的dB标度不同。

\paragraph{习题 7.11}
绳线密度 \(\mu\)，张力 \(F_T\)。
小振幅右行波
\(y=A\cos(kx-\omega t)\)。
小斜率下从左侧传向右侧的瞬时弹性功率
由横向张力分量与质元垂直速度的乘积给出：
\[
 P(x,t)=-F_T\frac{\partial y}{\partial x}
 \frac{\partial y}{\partial t}.
\]
代入
\(\partial_xy=-Ak\sin(kx-\omega t)\)，
\(\partial_ty=A\omega\sin(kx-\omega t)\)，
故
\[
 P=F_TA^2k\omega
 \sin^2(kx-\omega t).
\]
用弦波波速
\(u=\sqrt{F_T/\mu}\)，
且 \(\omega=ku\)：
\[
 \boxed{\overline P=
 \frac12F_TA^2k\omega
 =\frac12\mu uA^2\omega^2}.
\]
若绳有效截面积为 \(S\)，
平均强度 \(I=\overline P/S
 =\tfrac12\rho u\omega^2A^2\)，
与教材能流表达式一致。
正号表示波能沿波传播方向传输。

\paragraph{习题 7.12}
频率100 Hz、波速400 m/s，
波长
\(\lambda=u/f=4.0\,\mathrm m\)。
取直线AB长度30 m，
坐标 \(x\) 从A指向B；
两源原相差 \(\pi\)。
两路波在点 \(x\) 的路径差
\(\Delta r=(30-x)-x=30-2x\)。
增强条件为
\[
 k\Delta r+\pi=2n\pi
 \Longleftrightarrow
 \Delta r=(n-\tfrac12)\lambda.
\]
故
\(30-2x=(2n-1)2\)，
解出
\[
 \boxed{x=2,4,6,\ldots,28\,\mathrm m}.
\]
即两源之间（不包括源本身）
共有14个增强位置；
而 \(x=15\,\mathrm m\) 两程相同，
原初相相差 \(\pi\) 因而相消。
此结论假定传播中两波振幅相等。

\paragraph{习题 7.13}
驻波
\(y=0.020\cos(20x)\cos(750t)\)，
与两反向等振幅波叠加式
\(2a\cos(kx)\cos(\omega t)\) 比较：
\[
 \boxed{a=0.010\,\mathrm m,\quad
 k=20\,\mathrm{rad/m},\
 \omega=750\,\mathrm{rad/s}},\qquad
 \boxed{u=\omega/k=37.5\,\mathrm{m/s}}.
\]
节点满足
\(\cos(kx)=0\)，
相邻节点间距
\[
 \boxed{\Delta x=\frac\pi k
 =\frac\pi{20}\approx0.157\,\mathrm m}.
\]
质点速度
\[
 v_y=\frac{\partial y}{\partial t}
 =-15\cos(20x)\sin(750t).
\]
于
\(t=0.002\,\mathrm s,\ x=0.05\,\mathrm m\)
代入
\[
 \boxed{v_y=-15\cos(1.0)\sin(1.5)
 \approx-8.08\,\mathrm{m/s}}.
\]
符号表示此时质点向负 \(y\) 方向运动。

\paragraph{习题 7.14}
入射波沿 \(+x\) 传播，
可选初相使
\(y_i(0,0)=0\) 且
\((\partial_ty_i)(0,0)>0\)：
\[
 \boxed{y_i=A\sin(\omega t-kx)},
\quad
 \omega=2\pi\nu,\ k=2\pi/\lambda.
\]
在距原点
\(p=3\lambda/4\) 的反射面
从波疏到波密，位移反射相位改变 \(\pi\)。
从右向左的反射波为
\[
 y_r=-A\sin[\omega t+kx-2kp].
\]
因 \(2kp=3\pi\)，
\[
 \boxed{y_r=+A\sin(\omega t+kx)}.
\]
合振动
\[
 y_i+y_r
 =2A\sin(\omega t)\cos(kx).
\]
故节点位置满足
\(\cos(kx)=0\)，
在原点与反射面之间
\[
 \boxed{x=\lambda/4,\quad
 x=3\lambda/4=p},
\]
若继续考虑左侧空间，节点按
\(x=(2n+1)\lambda/4\) 周期排列。
反射面处的总位移确实为零，作为边界条件校验。

\paragraph{习题 7.15}
片厚
\(d=2.00\times10^{-3}\,\mathrm m\)。
厚度方向基频时
\(d=\lambda/2\)，
所以
\[
 \boxed{f=\frac{u}{2d}
 =\frac{5.74\times10^3}{4.00\times10^{-3}}
 =1.435\times10^6\,\mathrm{Hz}
 =1.435\,\mathrm{MHz}}.
\]
这是理想厚度剪切/纵向驻波的
几何估算，真实压电振动还受切型和电极负载影响。

\paragraph{习题 7.16}
采用声级定义
\(\beta=10\log_{10}(I/I_0)\,\mathrm{dB}\)
和题设
\(I_0=10^{-12}\,\mathrm{W/m^2}\)。
因此
\[
 \boxed{I_{115}=I_0\,10^{115/10}
 \approx0.316\,\mathrm{W/m^2}},\qquad
 \boxed{I_{141}=I_0\,10^{141/10}
 \approx1.26\times10^2\,\mathrm{W/m^2}}.
\]
141 dB对应强度约为
115 dB的 \(10^{(141-115)/10}\approx398\) 倍，
并不是声级数值的简单比值。

\paragraph{习题 7.17}
从图7.47按分钟刻度估读
两种地震波到时差
\(\Delta t=t_S-t_P\approx5\,\mathrm{min}
 =300\,\mathrm s\)。
震源距离为 \(r\) 时
\[
 t_P=\frac r{u_P},\quad
 t_S=\frac r{u_S},\quad
 \Delta t=r\left(\frac1{u_S}-\frac1{u_P}\right).
\]
用 \(u_P=8\,\mathrm{km/s}\)、
\(u_S=6\,\mathrm{km/s}\)，
\[
 \boxed{r=\frac{300}{1/6-1/8}
 \approx7.20\times10^3\,\mathrm{km}}.
\]
该估算对图上P与S起跳时刻的读数敏感，
将分钟错认成秒会导致60倍差错。
